% Generated draft: requires mathematical and source review.
\documentclass[11pt]{article}
\usepackage[margin=25mm]{geometry}
\usepackage{fontspec}
\setmainfont{Latin Modern Roman}
\usepackage{amsmath,amssymb,mathtools,unicode-math}
\setmathfont{Latin Modern Math}
\usepackage{enumitem,needspace,xurl,hyperref}
\setlength{\parindent}{0pt}
\setlength{\parskip}{5pt}
\setlength{\emergencystretch}{3em}
\newcommand{\N}{\mathbb{N}}
\newcommand{\Z}{\mathbb{Z}}
\newcommand{\Q}{\mathbb{Q}}
\newcommand{\R}{\mathbb{R}}
\newcommand{\C}{\mathbb{C}}
\newcommand{\F}{\mathbb{F}}
\newcommand{\A}{\mathbb{A}}
\newcommand{\PP}{\mathbb{P}}
\newcommand{\E}{\mathbb{E}}
\newcommand{\eps}{\varepsilon}
\newcommand{\dd}{\,\mathrm{d}}
\begin{document}
\section*{UnsolvedMath ID: 20001999}
\textbf{Source problem code:} catalogue entry
\subsection{Definitions and mathematical statement}
Margaret Symington's A.17 describes a tropical Calabi--Yau sphere as a singular integral-affine structure \(\mathcal A\) on \(B=S^2\setminus\{p_1,\ldots,p_{24}\}\). Around each singular point it requires (A) monodromy conjugate to the positive shear
\[
 T=\begin{pmatrix}1&1\\0&1\end{pmatrix},
\]
and (B) an injective local affine isomorphism \(\Phi:(U\setminus R,\mathcal A)\to(\R^2,\mathcal A_0)\), where \(R\) is a ray from the point. The matrix is the contextual reading of damaged extracted typography. The source asks for the geometry or topology of the set of such structures on \(S^2\), without specifying a topology on that set, markings, or an equivalence relation.

It reports the closed-surface possibilities under its local conditions: a torus or Klein bottle with no singularities, a sphere with 24, or \(\R\PP^2\) with 12. It says each occurs as the base of a singular Lagrangian fibration, whose singular fibers have genus-one Lefschetz nodal topology, described as spheres with one positive self-intersection. This wording concerns a positive node, not an asserted homological self-intersection number \(+1\). It attributes injectivity in (B) to a three-dimensional contact-geometric argument without giving its proof or citation.

The aid's reduction fixes substantially more data: label the 24 punctures, choose peripheral generators
\[
 \pi_1(B)=\langle\gamma_1,\ldots,\gamma_{24}\mid\gamma_1\cdots\gamma_{24}=1\rangle,
\]
and fix a positive factorization \(M_1\cdots M_{24}=I\) in \(SL(2,\Z)\), with each \(M_i\) conjugate to \(T\). Let \(\rho(\gamma_i)=M_i\) and \(V=\R^2\) with this action. An affine lift has the form
\[
 \widehat\rho(\gamma)(x)=\rho(\gamma)x+c(\gamma),\qquad
 c(\gamma\delta)=c(\gamma)+\rho(\gamma)c(\delta).
\]
The peripheral transformation has a fixed point exactly when
\(c(\gamma_i)\in\operatorname{im}(I-M_i)\). Modulo global translations, these lifts are the parabolic cohomology space
\[
 H^1_{\mathrm{par}}(B;V_\rho)=
 \ker\left(H^1(\pi_1(B);V_\rho)\longrightarrow
 \bigoplus_{i=1}^{24}H^1(\langle\gamma_i\rangle;V_\rho)\right).
\]
In coordinates \(b_i=c(\gamma_i)\), this is
\[
 \frac{\{(b_i): b_i\in\operatorname{im}(I-M_i),\quad
 \sum_{i=1}^{24}(M_1\cdots M_{i-1})b_i=0\}}
 {\{((I-M_i)v)_{i=1}^{24}:v\in\R^2\}}.
\]
The aid reports real dimension \(24-2-2=20\) for the fixed positive 24-node factorization. This is a translation-holonomy quotient with the linear representation fixed, not a quotient by all changes of marking, braid actions, or arbitrary linear affine conjugacy.

It further reports that this is only a holonomy shadow. If a node in affine coordinates has fixed-point parameter \(p\), its translational term is \((I-M)p\). Moving it to \(p+tv\) along an eigenline with \(Mv=v\) leaves this term unchanged. The aid interprets such nodal slides as geometric/developing-map data invisible to affine holonomy. Local fixed-point conditions alone do not establish global realizability or the injective cut-chart condition (B), so dimension 20 is not the dimension of the full unspecified structure space.

The original remark says that removing (B) but keeping monodromy condition (A) allows \(12k\) singularities on \(S^2\) for every integer \(k\ge2\). It also proposes more general tropical manifolds built from tropical Calabi--Yau strata with boundary and suitable compatibility, giving an affine cylinder with affine boundary as an example. It records Zharkov's question about tropical Gromov--Witten calculations on a Calabi--Yau and asks whether they extend to Lagrangian fibrations over these more general bases, specifically \(S^2\times T^2\) over the cylinder. No invariant, boundary conditions, fibration data, or compatibility axioms for that additional request are supplied. The aid reports modern polarized, triangulated, metric, and almost-toric moduli constructions, without claiming one unqualified moduli space answers all interpretations.
\subsection{Short English statement}
The question concerns the space of singular affine structures on a sphere with 24 positive nodes and injective cut charts. The aid identifies a 20-dimensional space of translation holonomies after fixing a labeled linear monodromy factorization, but node slides show that holonomy omits geometric data. The broader moduli and tropical Gromov--Witten requests remain underspecified.
\subsection{Sources}
\begin{itemize}
\item Local UnsolvedMath registry, permanent ID 20001999, exact source code catalogue entry; canonical \texttt{aim-geometry-notes.json}, \texttt{notes[336]}, tag \texttt{section}; AIM Workshop Problem Lists, Geometry, \emph{Amoebas and tropical geometry}, A.17, contributed by Margaret Symington. Created and updated 14 August 2026; statement status \texttt{reconstructed\_unverified}.
\item Supplied source URLs: \url{https://aimath.org/WWN/amoebas/amoebas.pdf} and \url{https://www.aimath.org/WWN/amoebas/articles/html/43a/}. The supplied verification reports checking both; neither was accessed here. The remark attributes a tropical Gromov--Witten question to Zharkov.
\item Attempt 1: \texttt{reduction}, problem status at run \texttt{partially\_solved}, rigor label \texttt{proved}, novelty confidence \texttt{low}; classification \texttt{PARTIAL-PROGRESS}, suggested difficulty L4, registry status \texttt{partially\_solved}. Candidate: parabolic-cohomology translation-holonomy reduction and invisible nodal-slide direction.
\item No complete authoritative references are supplied for the contact-geometric argument, surface classification and realizations, or the modern moduli constructions mentioned by the aid. All require source checking.
\item Ulam AI / UnsolvedMath Contributors; CC BY 4.0. Machine-generated research aid, not peer reviewed; independently verify before citation.
\end{itemize}
\end{document}
